% The Image Lives in Two Memories -- draft.
% Results are placeholders (\TBD) until phases 1-4 are merged from results/v0.2-main.
% Every number outside \TBD comes from the code, the model configs or docs/assumptions.md.
\documentclass[10pt,twocolumn]{article}
\usepackage[a4paper,margin=2cm,columnsep=0.7cm]{geometry}
\usepackage[T1]{fontenc}
\usepackage{amsmath,amsthm}   % before newtxmath (\openbox clash)
\usepackage{newtxtext,newtxmath}
\usepackage{booktabs,multirow,array}
\usepackage{graphicx}
\usepackage{xcolor}
\usepackage{url}
\usepackage[hidelinks]{hyperref}
\usepackage[numbers,sort&compress]{natbib}
\usepackage{caption}
\captionsetup{font=small}

\newcommand{\TBD}[1]{\textcolor{red}{[TBD: #1]}}
\newcommand{\R}{\mathbb{R}}
\newcommand{\NLL}{\mathrm{NLL}}
\newcommand{\ret}{\mathcal{R}}
\newcommand{\vis}{\mathcal{V}}
\newtheorem{remark}{Remark}

\title{\Large\bfseries The Image Lives in Two Memories:\\ Visual Information in Hybrid Linear-Attention Vision-Language Models}
\author{Minseok Shin\\
\small Horcrux Technologies, Seoul, Republic of Korea\\
\small \url{https://www.horcruxtechnologies.com}\\
\small \texttt{horcruxtech.ai@gmail.com}\quad\texttt{stevenshin16@korea.ac.kr}}
\date{}

\begin{document}
\twocolumn[
\maketitle
\begin{center}\begin{minipage}{0.86\textwidth}
\textbf{Abstract.}\;
Hybrid vision-language models that interleave linear-attention layers with softmax attention read an image into
two memories at once: the key-value (KV) rows of the visual tokens in the attention layers, and a fixed-size
recurrent state in the linear-attention layers that has absorbed every visual token. Visual-token pruning and
multimodal KV-cache compression assume that the image lives in the first of these. We ask where it actually lives.
On Qwen3.5 (Gated DeltaNet hybrids, 4B and 9B, early fusion) we alter one memory exactly once, right after the last
image token, and let the question and the answer read the altered memory. A dataset-level ratio of teacher-forced
answer losses, anchored at a no-image run and an unaltered run, measures how much of the image's contribution each
memory retains, with paired bootstrap intervals, on DocVQA, ChartQA, TextVQA and POPE.
\TBD{one sentence per finding: (i) KV zero vs state zero retained, per task family; (ii) bits needed by each memory
after reading, with and without error feedback; (iii) fewer visual tokens versus fewer bits at equal memory;
(iv) whether 9B agrees.}
\par\smallskip
\textbf{Keywords:} vision-language models; hybrid linear attention; Gated DeltaNet; KV cache; recurrent state;
visual token pruning; quantization; interpretability.
\end{minipage}\end{center}
\vspace{1em}
]

\section{Introduction}

A transformer vision-language model (VLM) keeps everything it has read about an image in one place: the keys and
values of the visual tokens, one row per token and layer. Visual-token pruning
\citep{chen2024fastv,zhang2025sparsevlm,shang2025prumerge,yang2025visionzip,xing2025pyramiddrop} and multimodal
KV-cache compression \citep{wan2024lookm,tu2025vlcache} are built on that picture: removing or compressing a visual
token's row removes or compresses what the model knows about that part of the image.

Hybrid models break the picture. Qwen3.5 \citep{qwen2026qwen35} repeats three Gated DeltaNet (GDN) layers
\citep{yang2025gdn} and one gated softmax-attention layer, and it is natively multimodal: visual features enter
at the input embedding and pass through the same stack as text (early fusion; Section~\ref{sec:arch}). When the
last image token has been processed, the image has been written into two memories. The attention layers hold the
KV rows of the visual tokens, as in a transformer. The GDN layers hold a fixed-size recurrent state that has been
updated by every visual token and is not indexed by token at all. In our setting this state is larger than the
visual KV rows of any image we feed (Section~\ref{sec:sizes}).

Which of the two memories does the model consult when it answers a question about the image? The answer decides
what token pruning and KV compression can safely do to a hybrid VLM, and whether the recurrent state, which a
serving engine stores at full precision regardless of the KV format \citep{vllm2026fp8kv}, is a place where image
information can be lost. It is not obvious a priori. The state is written by every visual token but has a finite
capacity and forgets through a learned decay; the KV rows keep every token exactly but sit in only one layer in
four.

We answer the question with training-free interventions. The prompt is read in two stages: the prefix and the
image first, then the question (Section~\ref{sec:twostage}). Between the stages we alter exactly one memory once:
we zero it, quantize it, or replace it by matched noise, or we remove visual tokens before they are read. The
question and the answer then read the altered memory. We measure the effect with the teacher-forced loss of the
reference answer, normalised between a run without the image and an unaltered run (Section~\ref{sec:metric}), so
that $1$ means the altered memory still supplies everything the image contributed and $0$ means nothing is left.

Two text-only studies of the same model family frame the question. Split-prefill experiments on Qwen3.5 and
Falcon-H1 keep only the KV cache or only the recurrent state of a prefilled context and find that exact retrieval
survives only through attention, while the output language and a persona survive through the recurrent state
\citep{afendulev2026recalls}. Quantizing both memories at every write found the recurrent state to be the fragile
one under long generation \citep{shin2026hybridquant}. Neither involves images. An image differs from a text
context in two ways that matter here: it is long (up to 1,280 tokens in our setting) and dense, and the questions
asked about it range from reading a single word (DocVQA, TextVQA) to judging whether an object is present at all
(POPE). The retrieval-versus-gist split of \citep{afendulev2026recalls} predicts that reading tasks depend on the
KV rows and presence judgements less so; our tasks test that prediction for visual content. For hybrid
Mamba--Transformer video models, token reduction has been found to behave like compression rather than dropping
because the state keeps information from discarded tokens \citep{jiang2026stateful}; we measure that state
contribution directly, for still images and Gated DeltaNet.

\paragraph{Contributions.}
\TBD{Final bullets after phases 1--4. Planned:}
\begin{itemize}\setlength\itemsep{1pt}
\item A one-shot intervention protocol for hybrid VLMs that separates the two memories at the end of the image
span, with a dataset-level retained ratio and paired bootstrap intervals (Sections~\ref{sec:twostage}--\ref{sec:metric}).
\item \TBD{Phase 1: which memory carries the image, per task (zero and noise interventions).}
\item \TBD{Phase 2: bits each memory needs after reading; error feedback during the image prefill; one-shot versus
every-write quantization.}
\item \TBD{Phase 3: visual-token deletion versus KV quantization at equal memory.}
\item \TBD{Phase 4: replication at 9B.}
\end{itemize}

\section{Related Work}

\paragraph{Visual-token reduction.}
FastV drops visual tokens after an early layer on the evidence of low attention to them \citep{chen2024fastv};
SparseVLM ranks visual tokens with text-guided attention \citep{zhang2025sparsevlm}; LLaVA-PruMerge and VisionZip
select and merge tokens at the vision encoder output \citep{shang2025prumerge,yang2025visionzip}; PyramidDrop
removes tokens progressively with depth \citep{xing2025pyramiddrop}; token merging goes back to ToMe
\citep{bolya2023tome}. All of these are designed and evaluated on transformer language backbones, where the
visual-token rows are the only memory of the image.
\TBD{verified citations from the literature check; add any hybrid/linear-VLM pruning work it finds.}

\paragraph{Multimodal KV-cache compression.}
LOOK-M and VL-Cache compress the KV cache of multimodal inputs with modality-aware eviction and budget allocation
\citep{wan2024lookm,tu2025vlcache}. KV quantization for text, e.g.\ KIVI \citep{liu2024kivi} and KVQuant
\citep{hooper2024kvquant}, applies unchanged to the visual rows we quantize.

\paragraph{Where visual information flows.}
Attention-knockout studies of transformer VLMs trace how image information reaches the answer position by
blocking attention between token groups in chosen layers
\citep{kaduri2025whats,zhang2025crossmodal,neo2025interpreting}. \TBD{one verified sentence per paper.} A
hybrid model adds a route these studies cannot see: the recurrent state carries information forward without
attention.

\paragraph{The two memories of hybrid models.}
Afendulev et al.\ separate the KV cache and the recurrent state of text contexts by keeping one and dropping the
other (the convolution cache is dropped with the state) and by swapping states between contexts
\citep{afendulev2026recalls}. We differ in the input (an image span inside the prompt, followed by the question),
in the interventions (graded: quantization, error feedback, matched noise and token deletion besides removal) and
in the measure (a loss ratio anchored at no-image and full runs). Hybrid video models have used the state as
evidence for which visual tokens to keep: as a carrier of discarded tokens \citep{jiang2026stateful} and, through
its change over a chunk, as a query-free eviction signal for the KV cache \citep{kwon2026deltas}.

\paragraph{Linear and hybrid VLMs; recurrent-state compression.}
Mamba-based VLMs replace the language backbone by a state-space model \citep{zhao2025cobra,qiao2024vlmamba};
there the state is the only memory of the image. Post-training quantization of the recurrent state of GDN and
related layers has been studied for text \citep{zhang2026damp,yao2026stepquant,pan2026leapquant}, and jointly with
the KV cache in \citep{shin2026hybridquant}. None of these isolates the image.

\section{Background}

\subsection{Qwen3.5 as a hybrid VLM}\label{sec:arch}
Qwen3.5 has $L$ decoder layers in repeating groups of three GDN layers and one gated softmax-attention layer, so
$L_a = L/4$ attention and $L_g = 3L/4$ GDN layers. For the 4B and 9B models $L=32$, $L_a=8$, $L_g=24$
(Table~\ref{tab:geometry}).

\paragraph{Early fusion.} A vision encoder with patch size $p=16$ and $m\times m$ spatial merging ($m=2$) turns an
image resized to $H\times W$ pixels (both multiples of $pm=32$) into
\begin{equation}
n_v = \frac{H}{p}\cdot\frac{W}{p}\cdot\frac{1}{m^2} = \frac{HW}{(pm)^2}
\label{eq:nv}
\end{equation}
visual tokens, which replace $n_v$ placeholder tokens in the input embedding. In the transformers implementation we
use, this scatter at the input embedding is the only place where visual features enter the language model; no
visual features are injected at intermediate layers (Appendix~\ref{app:premises}, premise 1). Visual tokens get
three-dimensional rotary positions (temporal, height, width; M-RoPE \citep{wang2024qwen2vl}) on the first quarter
of each attention head's dimensions. We cap the image at $HW \le 1{,}310{,}720$ pixels, so $n_v \le 1{,}280$.

\subsection{Gated DeltaNet}\label{sec:gdn}
Each GDN head keeps a matrix $S_t\in\R^{d_k\times d_v}$. With L2-normalised key $k_t\in\R^{d_k}$, query
$q_t\in\R^{d_k}$ (further scaled by $d_k^{-1/2}$), value $v_t\in\R^{d_v}$, decay $\alpha_t\in(0,1)$ and write
strength $\beta_t\in(0,1)$, computed per head from the input at $t$ as
$\alpha_t=\exp\!\big(-e^{A}\,\mathrm{softplus}(a_t+b)\big)$ and $\beta_t=\mathrm{sigmoid}(\beta'_t)$ with $a_t,\beta'_t$
linear in the input and $A,b$ learned, the state is first decayed and then corrected along $k_t$:
\begin{align}
S_t &= \alpha_t S_{t-1} + \beta_t\, k_t\big(v_t - \alpha_t S_{t-1}^{\top}k_t\big)^{\top} \nonumber\\
    &= \alpha_t\big(I-\beta_t k_t k_t^{\top}\big)S_{t-1} + \beta_t\, k_t v_t^{\top}, \label{eq:gdn}\\
o_t &= S_t^{\top} q_t. \label{eq:gdnout}
\end{align}
This is the order of operations of the reference recurrence in the implementation (decay, read
$(\alpha_tS_{t-1})^{\top}k_t$, write). Prefill computes the same recurrence in chunks, but the state is written to
the cache once per forward call, at its end (premise 5). Before the recurrence, $q,k,v$ pass through a causal
depthwise convolution of kernel size $\omega=4$ whose cache holds the last inputs; we return to it in
Section~\ref{sec:caveats}.

\subsection{Two memories and their sizes}\label{sec:sizes}
With $H_{kv}$ KV heads of dimension $d$ per attention layer and $H$ GDN value heads, a sequence of $T$ tokens
leaves
\begin{equation}
|\mathrm{KV}|(T) = 2L_aH_{kv}d\,T, \qquad |S| = L_gHd_kd_v
\label{eq:sizes}
\end{equation}
elements in the two memories (the factor 2 counts keys and values). The implementation stores the KV cache in the
model dtype (bf16, 16 bits) and the recurrent state in fp32 (32 bits; premise 3). For 4B and 9B,
$|\mathrm{KV}|(T)=16{,}384\,T$ and $|S|=24\cdot32\cdot128\cdot128=12{,}582{,}912$, i.e.\ 32.8\,kB of KV per token
and a 50.3\,MB state. The visual rows of an image therefore occupy $0.0328\,n_v$\,MB, at most 41.9\,MB at
$n_v=1{,}280$, and the two are equal only at
\begin{equation}
n_v^{\star} = \frac{32\,|S|}{16\cdot 2L_aH_{kv}d} = 1{,}536
\label{eq:crossover}
\end{equation}
visual tokens. Every image in this study leaves a state larger than its own KV rows. (The crossover of 768 tokens
in \citep{shin2026hybridquant} counts the state in bf16; the implementation used here keeps it in fp32.)

\begin{table}[t]
\centering\small
\caption{Cache geometry, from the model configurations. Sizes by \eqref{eq:sizes}; KV in bf16, state in fp32.}
\label{tab:geometry}
\setlength{\tabcolsep}{3.2pt}
\begin{tabular}{lcccccccc}
\toprule
Model & $L$ & $L_a$ & $L_g$ & $H_{kv}$ & $d$ & $H$ & $d_k{=}d_v$ & $|S|$ (MB) \\
\midrule
Qwen3.5-4B & 32 & 8 & 24 & 4 & 256 & 32 & 128 & 50.3 \\
Qwen3.5-9B & 32 & 8 & 24 & 4 & 256 & 32 & 128 & 50.3 \\
\bottomrule
\end{tabular}\\[2pt]
\footnotesize GDN key heads: 16; value heads $H=32$; the state is per value head. Vision encoder: patch 16,
merge 2; depth 24 (4B), 27 (9B).
\end{table}

\section{Method}

\subsection{Two-stage prefill and the intervention point}\label{sec:twostage}
A prompt is the token sequence $x = (P,\vis,Q)$: a prefix $P$ (chat template up to and including the
vision-start token), the image span $\vis=\{s,\dots,e-1\}$ of $n_v=e-s$ placeholder positions carrying the visual
features, and the question $Q$ (vision-end token, question, instruction and the assistant header). Let
$\mathcal{M}_e$ be the state of all caches after the first $e$ tokens:
\begin{equation}
\mathcal{M}_e = \big(\{K^{\ell}_{1:e},V^{\ell}_{1:e}\}_{\ell\in\mathcal{A}},\ \{S^{\ell}_e,\,C^{\ell}_e\}_{\ell\in\mathcal{G}}\big),
\end{equation}
with $\mathcal{A}$, $\mathcal{G}$ the attention and GDN layers and $C^{\ell}_e$ the convolution cache. Stage~1 runs
$(P,\vis)$ with unaltered computation. An intervention $\phi$ maps $\mathcal{M}_e$ to $\phi(\mathcal{M}_e)$. Stage~2
runs $Q$ on $\phi(\mathcal{M}_e)$ without further alteration; the reference answer is then scored by teacher forcing
on a copy of the resulting cache, and a greedy answer of at most 16 tokens is decoded on the original. Because
nothing is altered after the image, any effect of $\phi$ is an effect on what the image left behind. All
conditions, including the unaltered one, use the same two-stage path. \TBD{report the two-stage vs single-pass
check on 4B (Colab), not the local 0.8B check.}

\subsection{Interventions}\label{sec:interventions}
\paragraph{KV side.} For every attention layer, the rows of the visual positions are replaced,
$K^{\ell}_{\vis}\leftarrow f(K^{\ell}_{\vis})$ and $V^{\ell}_{\vis}\leftarrow f(V^{\ell}_{\vis})$, with
$f\in\{\mathrm{id},\,Q_4,\,N_4,\,0\}$. Rows of $P$ are not touched, and rows of $Q$ are written later by stage~2.
Keys are cached after the rotary embedding, so $f$ acts on rotated keys.

\paragraph{State side.} For every GDN layer, $S^{\ell}_e\leftarrow g(S^{\ell}_e)$ with
$g\in\{\mathrm{id},\,Q_8,\,Q_6,\,Q_8^{\mathrm{EF}},\,Q_6^{\mathrm{EF}},\,N_8,\,0\}$. The convolution cache
$C^{\ell}_e$ is never altered.

\paragraph{Quantizer.} $Q_b$ is asymmetric uniform round-to-nearest with one scale and zero point per group of
$G=128$ consecutive elements along the last dimension (the head dimension $d=256$ for keys and values, two groups
per head and token; $d_v$ for the state, one group per row $S[h,i,:]$). For a group $x\in\R^{G}$,
\begin{equation}
\begin{gathered}
\sigma = \frac{\max x-\min x}{2^b-1},\qquad z = \big\lfloor -\min x/\sigma \big\rceil,\\
Q_b(x) = \sigma\Big(\mathrm{clamp}\big(\lfloor x/\sigma\rceil + z,\,0,\,2^b-1\big) - z\Big),
\end{gathered}
\label{eq:quant}
\end{equation}
where $\lfloor\cdot\rceil$ rounds to the nearest integer and $\sigma$ is floored at $10^{-8}$. Quantization is
simulated in fp32 and cast back to the cache dtype.

\paragraph{Matched noise.} $N_b$ adds Gaussian noise with the energy of $b$-bit quantization error but none of its
structure. For a cache tensor $X$ (one layer's visual keys, values, or state),
\begin{equation}
N_b(X) = X + \rho_b(X)\,Z,\quad \rho_b(X) = \sqrt{\tfrac{1}{|X|}\textstyle\sum\big(Q_b(X)-X\big)^2},
\label{eq:noise}
\end{equation}
with $Z$ standard normal entries of the shape of $X$. The generator is re-seeded (seed 0) for every tensor, so
tensors of the same shape receive the same $Z$, across layers and samples, scaled by their own $\rho_b$. The RMS
$\rho_b$ is taken over the whole tensor, so noise is spread uniformly where quantization error is concentrated
in high-range groups.

\paragraph{Error feedback during the image.} Quantizing the state once after the image ($Q_b$ above) commits one
rounding error. With error feedback the image is read in chunks of at most 256 visual tokens, with boundaries
$s=e_0<e_1<\dots<e_C=e$, $C=\lceil n_v/256\rceil\le 5$. Let $\hat S_{e_u}$ be the state after running chunk $u$
from the stored state $\tilde S_{e_{u-1}}$, with $\tilde S_{e_0}=S_s$ the unaltered state after the prefix. Then
\begin{equation}
\tilde S_{e_u} = Q_b\big(\hat S_{e_u} + r_{u-1}\big),\quad r_u = \hat S_{e_u} + r_{u-1} - \tilde S_{e_u},
\label{eq:ef}
\end{equation}
with $r_0=0$; the final residual $r_C$ is discarded and $\tilde S_{e}$ is what the question reads. Each chunk is a
separate forward call, because the state is written once per call (premise 5). For $n_v\le 256$, $C=1$ and
$Q_b^{\mathrm{EF}}=Q_b$. Error feedback here therefore acts across at most five stores, unlike decoding in
\citep{shin2026hybridquant}, where it acts across every generated token.

\paragraph{Zeroing.} $f=0$ sets visual keys and values to zero; $g=0$ sets every GDN state to zero.
Section~\ref{sec:caveats} states precisely what each removes and what it leaves.

\paragraph{Visual-token deletion.} A keep set $I\subset\{0,\dots,n_v-1\}$ of size $n_v'=\max(1,\mathrm{round}(\kappa
n_v))$, $\kappa\in\{0.5,0.25\}$, is chosen uniformly, $I=\{\lfloor j\,n_v/n_v'\rfloor : j=0,\dots,n_v'-1\}$, or at random
without replacement (seed 0). The other visual tokens are removed from the input before the language model, so
neither memory ever sees them; the kept tokens keep their original M-RoPE positions, and the vision encoder still
sees the full image. Deletion shrinks the visual KV rows to $n_v'$ and leaves $|S|$ unchanged.

\paragraph{Every-write quantization.} For comparison with \citep{shin2026hybridquant} we also quantize every KV
append and every state write of the whole run, prefill and decoding, with the hooks of that work (scope
\emph{whole}); here the image and the text are quantized alike. \TBD{state the exact EF reset semantics of the
whole scope under the two-stage prefill.}

\subsection{Retained ratio}\label{sec:metric}
Let the reference answer have tokens $a_{1:m}$. For sample $i$ under condition $c$,
\begin{equation}
\NLL_i(c) = -\frac{1}{m}\sum_{j=1}^{m}\log p_c\big(a_j \mid x,\, a_{<j}\big),
\end{equation}
where $p_c$ reads the memory produced by condition $c$. Two anchors use the same prompt: \emph{full} (image,
no alteration) and \emph{blind} (no image: the image span and its start and end tokens are removed). With
$\bar N(c)=\frac1n\sum_i\NLL_i(c)$ over a task's $n$ samples, the retained ratio is
\begin{equation}
\ret(c) = \frac{\bar N(\mathrm{blind})-\bar N(c)}{\bar N(\mathrm{blind})-\bar N(\mathrm{full})}.
\label{eq:retained}
\end{equation}
$\ret=1$: the altered memory supplies the full reduction in answer loss that the image provides; $\ret=0$: none of
it; $\ret<0$: worse than no image. $\ret$ is a ratio of dataset means, not a mean of per-sample ratios, whose
denominators can be near zero. It is defined when $\bar N(\mathrm{blind})>\bar N(\mathrm{full})$, which we check
per task before reporting. The 95\% interval is a paired percentile bootstrap \citep{efron1993bootstrap}: 2,000
resamples of sample indices with replacement (seed 0), the same indices for the three means; resamples with a zero
denominator are dropped.

\paragraph{Reference answers and task scores.} The reference for the loss is the first answer (DocVQA), the most
frequent of the ten annotator answers (TextVQA), or the single answer (ChartQA, POPE). Greedy answers are scored
with ANLS at $\tau=0.5$ \citep{biten2019stvqa} (DocVQA), relaxed accuracy with 5\% relative tolerance
\citep{methani2020plotqa,masry2022chartqa} (ChartQA), VQA accuracy
$\frac{1}{10}\sum_{i=1}^{10}\min\!\big(1,\tfrac13\#\{j\neq i: a^{(j)}=\hat a\}\big)$ \citep{antol2015vqa} (TextVQA),
and yes/no accuracy and F1 with \emph{yes} as the positive class (POPE), all after one shared normalisation
(lower case; punctuation except decimal points and percent signs removed; articles removed).

\subsection{Memory accounting}\label{sec:budget}
For the equal-memory comparison we charge a condition
\begin{equation}
\mathrm{MB}(c) = \frac{b_{kv}\,|\mathrm{KV}|(n_v'+n_t) + b_s\,|S|}{8\cdot 10^{6}},
\label{eq:mb}
\end{equation}
with $n_v'$ the visual tokens kept, $n_t$ the text tokens, $b_{kv}\in\{16,4,0\}$ (full, 4-bit, zero) and
$b_s\in\{32,8,6,0\}$; noise is charged at its matched bit-width. This is deployment accounting: all KV rows are
charged at $b_{kv}$, although the one-shot intervention leaves the text rows at 16 bits. Scales and zero points
(about $32/G=0.25$ bits per element for fp16 pairs) and the convolution cache are not counted.

\subsection{Caveats of the interventions}\label{sec:caveats}
\paragraph{Zeroed keys still take attention.} A zero key gives the logit $q^{\top}0/\sqrt{d}=0$, not $-\infty$, so
each zeroed visual position keeps a softmax weight $w_i=1/\mathcal{Z}$, with $\mathcal{Z}$ the softmax normaliser of the
query, and contributes the zero value. The attention output becomes $\sum_{i\notin\vis}w_iv_i$ with
$\sum_{i\notin\vis}w_i = 1-n_v/\mathcal{Z}<1$: KV zero removes the visual
content and also shrinks the attention output by the mass the zeroed rows absorb. Masking the visual positions
would remove the content and renormalise; \citep{afendulev2026recalls} avoid zero keys for this reason.
\TBD{decide whether to add a KV-drop condition (visual rows removed from the cache); until then KV zero is read as
removal plus attention dilution.}

\paragraph{The convolution cache leaks the last visual tokens.} Unlike \citep{afendulev2026recalls}, who drop the
convolution cache together with the state, we leave it in place, following \citep{shin2026hybridquant}. $C^{\ell}_e$ holds the inputs of the last positions
of the span, and the causal convolution of kernel $\omega=4$ mixes them into the GDN inputs of the first $\omega-1=3$
tokens after the span. State zero therefore leaves a short local channel from the last three visual tokens into the
first question-side tokens (and, through their hidden states, into everything above them).

\paragraph{State zero also forgets the prefix.} $S^{\ell}_e$ summarises $P$ as well as $\vis$; zeroing it removes
both. The prefix is the fixed chat-template header, the same for every sample.

\section{Experimental Setup}\label{sec:setup}
\paragraph{Models.} Qwen3.5-4B (phases 1--3) and Qwen3.5-9B (phase 4), bf16 weights, the post-trained (non-base)
checkpoints. GDN runs in the library's PyTorch kernels. \TBD{confirm from the Colab environment that no fla kernel
path was taken, and report torch/transformers versions recorded in the result rows.}

\paragraph{Data.} Fixed random subsets (seed 0) of DocVQA validation (200) \citep{mathew2021docvqa}, ChartQA test
(200) \citep{masry2022chartqa}, TextVQA validation (200) \citep{singh2019textvqa} and POPE (300) \citep{li2023pope},
from the lmms-lab copies used by lmms-eval \citep{zhang2025lmmseval} at pinned revisions; the sample ids are released.
Every question gets the same instruction, ``Answer the question using a single word or phrase.'', in the default
chat template without a system prompt; the template closes an empty thinking block before the answer
(non-thinking mode).

\paragraph{Conditions.} Phase 1 (which memory): $f\in\{\mathrm{id},Q_4,N_4,0\}$ with the state unaltered;
$g\in\{Q_8,Q_8^{\mathrm{EF}},Q_6,Q_6^{\mathrm{EF}},N_8,0\}$ with the KV unaltered; and both zeroed. Phase 2 (bits):
the state levels alone and with $f=Q_4$, and the two every-write cells (KV4 with $Q_8$ and $Q_8^{\mathrm{EF}}$).
Phase 3 (tokens versus bits): $\kappa\in\{1,0.5,0.25\}$ uniform, with KV at 16 or 4 bits, and $\kappa=0.5$ random.
Phase 4 (9B): full, KV zero, state zero, KV4 with $Q_8^{\mathrm{EF}}$, and $\kappa=0.5$ uniform. Each condition
also has the two anchors per task. Batch size 1 and greedy decoding throughout.

\paragraph{Compute.} One NVIDIA A100-SXM4-80GB on Google Colab. A 4B preflight (22 conditions $\times$ 4 tasks
$\times$ 5 samples) took 1.47\,s per sample on average, with a peak of 9.3\,GB of GPU memory. \TBD{total GPU hours
and compute units per phase.}

\section{Results}
\subsection{Which memory carries the image}\label{sec:res-probe}
\TBD{Figure: retained per condition and task with 95\% CIs (paper/figures/retained\_Qwen3.5-4B.pdf). Table:
phase-1 retained. Text: KV zero vs state zero; KV noise vs KV4; state noise vs state 8 (structure vs energy, cf.\
the noise control in \citep{shin2026hybridquant}); both zeroed versus blind (should be near 0 unless the conv leak
carries information).}

\subsection{How many bits each memory needs after reading}\label{sec:res-bits}
\TBD{State 8/6 with and without EF, alone and with KV4. EF only differs from RTN for $n_v>256$: report $n_v$
distribution per task and split by $C$. One-shot vs every-write.}

\subsection{Fewer visual tokens or fewer bits}\label{sec:res-tokens}
\TBD{Table at matched MB by \eqref{eq:mb}: $\kappa=0.5$ at 16 bits vs $\kappa=1$ at 4 bits (KV rows of the image
$0.5\times16$ vs $1\times4$ bits per element). Uniform vs random.}

\subsection{Scale}\label{sec:res-scale}
\TBD{9B vs 4B on the five phase-4 conditions.}

\subsection{Task scores}\label{sec:res-scores}
\TBD{ANLS / relaxed accuracy / VQA accuracy / POPE F1 per condition; does the ordering agree with $\ret$?}

\section{Discussion and Limitations}
\TBD{Practical reading for pruning and KV compression of hybrid VLMs.}
\paragraph{Limitations.} (i) KV zero keeps softmax mass on the zeroed rows and state zero leaves the convolution
cache (Section~\ref{sec:caveats}); neither is a perfect deletion. (ii) Each intervention acts once, after the
image; it does not model a cache that stays quantized while the answer is generated, except for the two
every-write cells. (iii) $\ret$ is measured with the loss of one reference answer per sample; the task scores
check that it tracks answer quality, but a reference that the model would phrase differently inflates both anchors.
(iv) Subsets of 200--300 samples per task, one model family at two sizes, one image resolution cap.
(v) Quantization is simulated; no memory or speed is saved in these runs.

\section{Conclusion}
\TBD{after results.}

\section*{Acknowledgments}
This work was carried out by a single author without institutional funding; compute was self-funded on Google
Colab. The author used Claude (Anthropic) as an assistant for code development and manuscript editing; all
experiments, analyses and claims were designed, run and verified by the author.

\section*{Reproducibility and Data Availability}
Code, configurations, subset ids, raw per-sample results and the scripts that regenerate every table and figure
are at \url{https://github.com/skytinstone/twomemories} \TBD{make public at submission}. Each phase is run by
\texttt{python scripts/run\_phase.py --phase <name>}, the tables by \texttt{python analysis/make\_tables.py} and the
figures by \texttt{python analysis/figures.py}. Quantization uses the released \texttt{hybridquant} package
\citep{shin2026hybridquant} unchanged. All datasets are public.

\bibliographystyle{unsrtnat}
\bibliography{refs}

\appendix
\onecolumn
\section{Implementation premises}\label{app:premises}
Checked against transformers 5.22.2 (Qwen3.5 modeling and cache code) before any run; file and line references
are in \texttt{docs/assumptions.md} of the repository.
\begin{center}\small
\begin{tabular}{p{0.04\textwidth}p{0.42\textwidth}p{0.44\textwidth}}
\toprule
& Premise & Consequence for the method \\
\midrule
1 & Visual features are scattered into the input embedding only; no intermediate-layer injection. &
The caches after the span (KV, state, convolution cache) are the only carriers of the image. \\
2 & The image is one contiguous span of placeholder tokens of length $\prod(\text{grid})/m^2$. &
$\vis$ is known exactly; \eqref{eq:nv}. \\
3 & The recurrent state is stored in fp32 (PyTorch path). & State charged at 32 bits; \eqref{eq:crossover}. \TBD{GPU path}. \\
4 & The cache is a \texttt{DynamicCache} with attention and linear-attention layers; the write methods hooked by
\texttt{hybridquant} exist unchanged. & No adapter; every-write cells use the hooks as released. \\
5 & Prefill computes GDN in 64-token chunks internally but writes the state once per forward call. &
Error feedback during the image needs one forward call per 256-token chunk \eqref{eq:ef}. \\
\bottomrule
\end{tabular}
\end{center}

\section{Full results}\label{app:results}
\TBD{generated by analysis/make\_tables.py from results/**/*.jsonl.}

\end{document}
